\section{Resultados e Discussão}
\label{sec:resultados}

\subsection{População e cobertura efetiva}

A filtragem da seleção original reduziu 38.456 registros a 540 candidatos
criados até 2017: 20 antigos e 520 criados em 2014--2017. Todos foram
percorridos, sem substituição. Foram registradas as 1.640 posições previstas,
com 1.460 snapshots recuperados. Houve 174 posições sem commit dentro da era
e seis falhas de recuperação, correspondentes a dois repositórios indisponíveis
em suas três eras previstas.

Entre os snapshots recuperados, 1.441 tiveram extração CK com estado completo
no diagnóstico adotado; 15 não continham fontes Java no escopo, três
apresentaram erro do extrator e um teve cobertura parcial. Extração completa
não implica elegibilidade: confirmação do framework, tamanho e demais
critérios reduziram a comparação principal a 584 versões, pertencentes a
245 projetos distintos. Os totais por era são versões de projetos que podem
se repetir; não devem ser somados como repositórios únicos.

% Generated from the independently audited original-selection panel.
\begin{table}[H]
\centering\small
\caption{Cobertura das posições projeto--era na seleção original.}
\begin{tabular}{lrrrr}
\hline
Era & Previstas & Com SHA & Elegíveis & Sem commit \\ \hline
E1 & 20 & 20 & 6 & 0 \\
E2 & 540 & 525 & 162 & 13 \\
E3 & 540 & 472 & 210 & 66 \\
E4 & 540 & 443 & 206 & 95 \\
\hline
\end{tabular}
\end{table}

\begin{table}[H]
\centering\small
\caption{Médias entre projetos elegíveis por era; a composição varia entre linhas.}
\begin{tabular}{lrrrrrrrr}
\hline
Era & Projetos & CBO & RFC & WMC & LOC/método & DIT & NOC & LCOM* \\ \hline
E1 & 6 & 5,103 & 8,460 & 8,415 & 6,455 & 1,373 & 0,091 & 0,274 \\
E2 & 162 & 5,773 & 8,285 & 7,416 & 5,847 & 1,291 & 0,084 & 0,276 \\
E3 & 210 & 5,956 & 8,779 & 7,024 & 5,986 & 1,302 & 0,092 & 0,264 \\
E4 & 206 & 6,150 & 9,663 & 7,017 & 6,225 & 1,297 & 0,088 & 0,257 \\
\hline
\end{tabular}
\end{table}


Os 20 antigos tiveram suas quatro eras verificadas. Foram elegíveis seis
versões em E1, nove em E2, onze em E3 e dez em E4. A versão de E1 do próprio
Spring Boot foi identificada como código do framework; a presença preliminar
de Boot no guia de serviço REST também foi registrada, mas seu tamanho não
atingiu os limites principais. Assim, não se presumiu ausência de Boot antes
do lançamento estável nem se confundiu código do framework com uma aplicação.

As falhas permanecem individualizadas nos dados. Dois erros ocorreram em
versões de um projeto com arquivos de template e outro em um exemplo Java
não processado pelo visitante do CK. A cobertura parcial envolveu caminhos
Java que diferem apenas por maiúsculas/minúsculas. Nenhum desses casos foi
convertido em métrica zero ou trocado por uma versão anterior conveniente.

\subsection{Mudanças nos mesmos projetos}

Cinco projetos antigos foram elegíveis em E1 e E4. Nesse conjunto, CBO e RFC
aumentaram em média, enquanto WMC diminuiu. LOC por método, NOC e LCOM*
variaram pouco em termos absolutos. Esses resultados descrevem dimensões
distintas; não sustentam uma direção única de melhoria estrutural.

Nos 123 projetos criados em 2014--2017 elegíveis em E2 e E4, as médias de CBO,
RFC, WMC e LOC por método aumentaram. DIT e NOC variaram pouco, e LCOM*
diminuiu. A redução desse indicador de falta de coesão coexistiu com aumento
de acoplamento e complexidade. Não se agrega esse comportamento em uma nota
única de maturidade.

% Tabelas longitudinais estritas: apenas repositórios válidos com CK em todas as eras de seu grupo (17 pioneiros em E1-E4; 453 principais em E2-E4)

\begin{table}[H]
\centering\small
\caption{Grupo pioneiro (criados até 2013, $n = 17$): médias das métricas CK dos mesmos 17 repositórios válidos em E1, E2, E3 e E4.}
\label{tab:valores_pioneiro}
\begin{tabular}{lrrrr}
\hline
\textbf{Métrica} & \textbf{E1} & \textbf{E2} & \textbf{E3} & \textbf{E4} \\ \hline
CBO & 4,878 & 5,394 & 5,626 & 5,803 \\
RFC & 7,639 & 8,387 & 8,972 & 9,956 \\
WMC & 8,094 & 8,076 & 7,936 & 7,201 \\
LOC/método & 6,316 & 5,957 & 6,208 & 6,446 \\
DIT & 1,416 & 1,355 & 1,376 & 1,341 \\
NOC & 0,083 & 0,068 & 0,104 & 0,081 \\
LCOM* & 0,298 & 0,276 & 0,277 & 0,249 \\
\hline
\end{tabular}
\end{table}

\begin{table}[H]
\centering\small
\caption{Grupo pioneiro (criados até 2013, $n = 17$): diferenças das médias de E1 em relação às eras subsequentes nos mesmos 17 repositórios válidos.}
\label{tab:diferencas_pioneiro}
\begin{tabular}{lrrr}
\hline
\textbf{Métrica} & \textbf{(E2 $-$ E1)} & \textbf{(E3 $-$ E1)} & \textbf{(E4 $-$ E1)} \\ \hline
CBO & +0,516 & +0,748 & +0,925 \\
RFC & +0,748 & +1,333 & +2,317 \\
WMC & -0,019 & -0,158 & -0,893 \\
LOC/método & -0,359 & -0,108 & +0,130 \\
DIT & -0,062 & -0,040 & -0,076 \\
NOC & -0,015 & +0,021 & -0,002 \\
LCOM* & -0,022 & -0,021 & -0,049 \\
\hline
\end{tabular}
\end{table}

\begin{table}[H]
\centering\small
\caption{Grupo principal (criados entre 2014 e 2017, $n = 453$): médias das métricas CK dos mesmos 453 repositórios válidos em E2, E3 e E4.}
\label{tab:valores_principal}
\begin{tabular}{lrrr}
\hline
\textbf{Métrica} & \textbf{E2} & \textbf{E3} & \textbf{E4} \\ \hline
CBO & 5,405 & 5,591 & 5,671 \\
RFC & 6,603 & 7,224 & 7,587 \\
WMC & 5,819 & 5,801 & 5,759 \\
LOC/método & 5,308 & 5,489 & 5,616 \\
DIT & 1,221 & 1,229 & 1,226 \\
NOC & 0,059 & 0,065 & 0,063 \\
LCOM* & 0,245 & 0,243 & 0,235 \\
\hline
\end{tabular}
\end{table}

\begin{table}[H]
\centering\small
\caption{Grupo principal (criados entre 2014 e 2017, $n = 453$): diferenças das médias entre E2, E3 e E4 nos mesmos 453 repositórios válidos.}
\label{tab:diferencas_principal}
\begin{tabular}{lrrr}
\hline
\textbf{Métrica} & \textbf{(E3 $-$ E2)} & \textbf{(E4 $-$ E3)} & \textbf{(E4 $-$ E2)} \\ \hline
CBO & +0,186 & +0,080 & +0,266 \\
RFC & +0,622 & +0,363 & +0,984 \\
WMC & -0,018 & -0,042 & -0,060 \\
LOC/método & +0,181 & +0,127 & +0,307 \\
DIT & +0,008 & -0,003 & +0,005 \\
NOC & +0,006 & -0,002 & +0,004 \\
LCOM* & -0,002 & -0,007 & -0,009 \\
\hline
\end{tabular}
\end{table}


Os gráficos do relatório mantêm composição constante em todas as eras de
cada coorte: quatro projetos antigos em E1--E4 e 118 posteriores em E2--E4.
Esses conjuntos são menores que os pares das tabelas, pois exigem elegibilidade
também nas eras intermediárias. A cobertura pode alterar a leitura: nos
quatro antigos do painel balanceado, LCOM* passou de 0,279 para 0,302,
enquanto nos cinco pares E1--E4 houve pequena redução de 0,290 para 0,287.
Portanto, explicitar os IDs compartilhados é necessário mesmo em uma análise
descritiva.

\subsection{Presença de Spring Boot e evidências de IA}

Entre endpoints adjacentes elegíveis, a passagem observável de Spring
tradicional para presença declarada de Boot ocorreu em quatro pares E1--E2,
nove E2--E3 e quatro E3--E4. A diferença média de CBO nesses conjuntos foi,
respectivamente, $-0,260$, $+1,136$ e $-0,498$. Os conjuntos e períodos
distintos não apresentaram uma direção uniforme. A presença de Boot em um
módulo ou exemplo não equivale à migração integral do projeto; esses números
não isolam um efeito do framework.

Em E4, 94 candidatos apresentaram evidência detectada de IA, 349 não
apresentaram os sinais procurados e 97 permaneceram indeterminados: 95 sem
snapshot na era e dois sem recuperação do repositório. Entre os 206
snapshots elegíveis para CK, 49 tinham evidência e 157 não tinham evidência
detectada. Não foram encontrados sinais nas versões de E1--E3 pelo detector
adotado; isso não comprova inexistência de assistência nesses períodos.

\begin{table}[htbp]
\centering\footnotesize
\caption{E4: médias entre os repositórios válidos finais ($N = 470$) com evidência detectada de IA (A, $n_a=75$) e sem evidência detectada (B, $n_b=395$).}
\label{tab:ia_e4}
\begin{tabular}{lrrr}
\hline
\textbf{Métrica} & \textbf{Média A} & \textbf{Média B} & \textbf{A $-$ B} \\ \hline
CBO & 6,101 & 5,595 & $+$0,506 \\
RFC & 9,840 & 7,261 & $+$2,579 \\
WMC & 6,960 & 5,593 & $+$1,366 \\
LOC/método & 6,182 & 5,544 & $+$0,639 \\
DIT & 1,301 & 1,216 & $+$0,085 \\
NOC & 0,079 & 0,061 & $+$0,018 \\
LCOM* & 0,251 & 0,233 & $+$0,018 \\
\hline
\end{tabular}
\par\vspace{2pt}{\footnotesize Fonte: Elaborada pelos autores (2026). Nota: diferenças calculadas antes do arredondamento.}
\end{table}


As sete médias de E4 foram maiores no grupo com evidência. Entretanto, esse
grupo também apresentou maior porte: média de 125.578 LOC totais por projeto,
contra 15.222 no grupo sem evidência detectada. A comparação não controla
tamanho, domínio, arquitetura ou equipe e não permite concluir que a IA
aumentou complexidade ou piorou qualidade.

O momento observado também difere. Os snapshots elegíveis com evidência estão,
em média, 34,4 dias antes do corte fixado; os sem evidência, 475,0 dias antes.
Os primeiros variam de novembro de 2025 a outubro de 2026, enquanto os segundos
incluem versões desde fevereiro de 2022. Compartilhar a era E4 não iguala a
oportunidade temporal de observar vestígios dos agentes. Essa diferença deve
acompanhar a interpretação, sem reclassificar versões antigas por sinais futuros.

Nos 43 pares elegíveis com ausência de vestígio em E3 e evidência observada em
E4, CBO médio passou de 6,356 para 6,618. Nos 139 pares sem vestígio nos dois
endpoints, passou de 5,925 para 5,975. A diferença entre esses conjuntos já
aparecia na versão de E3. Além disso, a primeira evidência observada não data
o primeiro uso real. A análise representa associação temporal, sem identificação
causal.

\subsection{Alcance das evidências}

O resultado central é uma descrição longitudinal rastreável, com cobertura
explícita e mudanças heterogêneas entre dimensões. A seleção original por
tópicos atuais, atividade e arquivamento preserva viés de sobrevivência.
Configurações não resolvidas pelo detector, perdas de histórico, falhas do CK
e assistência sem vestígios também restringem o alcance. Os poucos projetos
antigos comparáveis impedem generalizar sua trajetória para toda a comunidade
Spring. Os resultados de populações anteriores não foram incorporados a
essas tabelas.

A auditoria posterior identificou limitações adicionais de composição, sem
alterar a amostra ou as medições. Dos 91 candidatos sinalizados para revisão
educacional, 29 tiveram alguma versão elegível; a confirmação estática de Spring
não determina se o projeto é uma aplicação, biblioteca ou demonstração.
Em 451 dos 584 snapshots elegíveis há classes sem valor finito de LCOM*,
excluídas apenas do denominador dessa média. Em 98 versões, LOC classificados
como produção ficam abaixo de 2.000, embora produção e testes juntos atendam
ao limite adotado. A classificação por caminho pode errar e não foi validada
individualmente em todos os módulos.

Além disso, 105 snapshots elegíveis não possuem versão declarada preenchida.
O estudo atual registra porte e framework, mas não executa o controle de
tamanho e versão originalmente previsto para os contrastes inferenciais.
Dispersão descritiva está disponível no suplemento de auditoria; médias e
extremos não eliminam confundimento. Uma análise de sensibilidade com regras
fixadas previamente pode complementar o estudo, sem substituir os resultados
principais ou retirar casos em função das métricas observadas.
